% Options for packages loaded elsewhere
\PassOptionsToPackage{unicode}{hyperref}
\PassOptionsToPackage{hyphens}{url}
%
\documentclass[
]{article}
\usepackage{amsmath,amssymb}
\usepackage{iftex}
\ifPDFTeX
  \usepackage[T1]{fontenc}
  \usepackage[utf8]{inputenc}
  \usepackage{textcomp} % provide euro and other symbols
\else % if luatex or xetex
  \usepackage{unicode-math} % this also loads fontspec
  \defaultfontfeatures{Scale=MatchLowercase}
  \defaultfontfeatures[\rmfamily]{Ligatures=TeX,Scale=1}
\fi
\usepackage{lmodern}
\ifPDFTeX\else
  % xetex/luatex font selection
\fi
% Use upquote if available, for straight quotes in verbatim environments
\IfFileExists{upquote.sty}{\usepackage{upquote}}{}
\IfFileExists{microtype.sty}{% use microtype if available
  \usepackage[]{microtype}
  \UseMicrotypeSet[protrusion]{basicmath} % disable protrusion for tt fonts
}{}
\makeatletter
\@ifundefined{KOMAClassName}{% if non-KOMA class
  \IfFileExists{parskip.sty}{%
    \usepackage{parskip}
  }{% else
    \setlength{\parindent}{0pt}
    \setlength{\parskip}{6pt plus 2pt minus 1pt}}
}{% if KOMA class
  \KOMAoptions{parskip=half}}
\makeatother
\usepackage{xcolor}
\usepackage[margin=1in]{geometry}
\usepackage{color}
\usepackage{fancyvrb}
\newcommand{\VerbBar}{|}
\newcommand{\VERB}{\Verb[commandchars=\\\{\}]}
\DefineVerbatimEnvironment{Highlighting}{Verbatim}{commandchars=\\\{\}}
% Add ',fontsize=\small' for more characters per line
\usepackage{framed}
\definecolor{shadecolor}{RGB}{248,248,248}
\newenvironment{Shaded}{\begin{snugshade}}{\end{snugshade}}
\newcommand{\AlertTok}[1]{\textcolor[rgb]{0.94,0.16,0.16}{#1}}
\newcommand{\AnnotationTok}[1]{\textcolor[rgb]{0.56,0.35,0.01}{\textbf{\textit{#1}}}}
\newcommand{\AttributeTok}[1]{\textcolor[rgb]{0.13,0.29,0.53}{#1}}
\newcommand{\BaseNTok}[1]{\textcolor[rgb]{0.00,0.00,0.81}{#1}}
\newcommand{\BuiltInTok}[1]{#1}
\newcommand{\CharTok}[1]{\textcolor[rgb]{0.31,0.60,0.02}{#1}}
\newcommand{\CommentTok}[1]{\textcolor[rgb]{0.56,0.35,0.01}{\textit{#1}}}
\newcommand{\CommentVarTok}[1]{\textcolor[rgb]{0.56,0.35,0.01}{\textbf{\textit{#1}}}}
\newcommand{\ConstantTok}[1]{\textcolor[rgb]{0.56,0.35,0.01}{#1}}
\newcommand{\ControlFlowTok}[1]{\textcolor[rgb]{0.13,0.29,0.53}{\textbf{#1}}}
\newcommand{\DataTypeTok}[1]{\textcolor[rgb]{0.13,0.29,0.53}{#1}}
\newcommand{\DecValTok}[1]{\textcolor[rgb]{0.00,0.00,0.81}{#1}}
\newcommand{\DocumentationTok}[1]{\textcolor[rgb]{0.56,0.35,0.01}{\textbf{\textit{#1}}}}
\newcommand{\ErrorTok}[1]{\textcolor[rgb]{0.64,0.00,0.00}{\textbf{#1}}}
\newcommand{\ExtensionTok}[1]{#1}
\newcommand{\FloatTok}[1]{\textcolor[rgb]{0.00,0.00,0.81}{#1}}
\newcommand{\FunctionTok}[1]{\textcolor[rgb]{0.13,0.29,0.53}{\textbf{#1}}}
\newcommand{\ImportTok}[1]{#1}
\newcommand{\InformationTok}[1]{\textcolor[rgb]{0.56,0.35,0.01}{\textbf{\textit{#1}}}}
\newcommand{\KeywordTok}[1]{\textcolor[rgb]{0.13,0.29,0.53}{\textbf{#1}}}
\newcommand{\NormalTok}[1]{#1}
\newcommand{\OperatorTok}[1]{\textcolor[rgb]{0.81,0.36,0.00}{\textbf{#1}}}
\newcommand{\OtherTok}[1]{\textcolor[rgb]{0.56,0.35,0.01}{#1}}
\newcommand{\PreprocessorTok}[1]{\textcolor[rgb]{0.56,0.35,0.01}{\textit{#1}}}
\newcommand{\RegionMarkerTok}[1]{#1}
\newcommand{\SpecialCharTok}[1]{\textcolor[rgb]{0.81,0.36,0.00}{\textbf{#1}}}
\newcommand{\SpecialStringTok}[1]{\textcolor[rgb]{0.31,0.60,0.02}{#1}}
\newcommand{\StringTok}[1]{\textcolor[rgb]{0.31,0.60,0.02}{#1}}
\newcommand{\VariableTok}[1]{\textcolor[rgb]{0.00,0.00,0.00}{#1}}
\newcommand{\VerbatimStringTok}[1]{\textcolor[rgb]{0.31,0.60,0.02}{#1}}
\newcommand{\WarningTok}[1]{\textcolor[rgb]{0.56,0.35,0.01}{\textbf{\textit{#1}}}}
\usepackage{graphicx}
\makeatletter
\def\maxwidth{\ifdim\Gin@nat@width>\linewidth\linewidth\else\Gin@nat@width\fi}
\def\maxheight{\ifdim\Gin@nat@height>\textheight\textheight\else\Gin@nat@height\fi}
\makeatother
% Scale images if necessary, so that they will not overflow the page
% margins by default, and it is still possible to overwrite the defaults
% using explicit options in \includegraphics[width, height, ...]{}
\setkeys{Gin}{width=\maxwidth,height=\maxheight,keepaspectratio}
% Set default figure placement to htbp
\makeatletter
\def\fps@figure{htbp}
\makeatother
\setlength{\emergencystretch}{3em} % prevent overfull lines
\providecommand{\tightlist}{%
  \setlength{\itemsep}{0pt}\setlength{\parskip}{0pt}}
\setcounter{secnumdepth}{-\maxdimen} % remove section numbering
\ifLuaTeX
  \usepackage{selnolig}  % disable illegal ligatures
\fi
\usepackage{bookmark}
\IfFileExists{xurl.sty}{\usepackage{xurl}}{} % add URL line breaks if available
\urlstyle{same}
\hypersetup{
  pdftitle={BonusQuestion},
  pdfauthor={GB},
  hidelinks,
  pdfcreator={LaTeX via pandoc}}

\title{BonusQuestion}
\author{GB}
\date{2026-04-09}

\begin{document}
\maketitle

\section{1 notes on the recap
material}\label{notes-on-the-recap-material}

These questions focus on material that students commonly struggle with
(data wrangling, plotting, modeling, or debugging).

Ans: Data wrangling (dplyr) Select () chooses columns Filter() keeps
rows based on conditions Mutate() create new variables Summarise ()
reduces data to summary statistics Group\_by() groups data before
summarizing Always check column names and data types before wrangling

Examples; library(dplyr)

data\_summary \textless- df \%\textgreater\% filter(year == 2025)
\%\textgreater\% group\_by(treatment) \%\textgreater\%
summarise(mean\_yield = mean(yield, na.rm = TRUE))

Common mistakes: summarise without group\_by wrong column names factors
vs numeric mismatch

Plotting (ggplot2)

Basic structure:

ggplot(data, aes(x, y)) + geom\_point()

Common layers:

geom\_point() scatter plot geom\_line() line plot
geom\_bar(stat=``identity'') bar chart
facet\_wrap(\textasciitilde variable) multiple plots theme\_classic()
clean theme

Example:

ggplot(df, aes(treatment, yield)) + geom\_boxplot() + theme\_classic()

Common mistakes:

mapping inside vs outside aes() forgetting + wrong variable names
plotting summarized data incorrectly Linear modeling

Basic model:

model \textless- lm(yield \textasciitilde{} treatment, data = df)
summary(model)

Important outputs:

Estimate = effect size p-value = significance R² = model fit Residuals =
model error

Prediction:

predict(model)

Common mistakes:

response and predictor reversed categorical variables not converted to
factor missing values in model data Debugging

Steps:

read the error message carefully run code line-by-line check object
names check parentheses and commas use View(df) or str(df) print
intermediate outputs

Example:

str(df) names(df)

\subsubsection{Pipes \%\textgreater\%}\label{pipes}

Confusing library(tidyverse) microbiome.fungi \textless-
read.csv(file.choose()) str(microbiome.fungi) microbiome.fungi2
\textless- select(microbiome.fungi, SampleID, Crop,
Compartment:Fungicide, richness) str(microbiome.fungi2)
head(filter(microbiome.fungi2, Treatment == ``Conv.'')) microbiome.fungi
\%\textgreater\% select(SampleID, Crop, Compartment:Fungicide, richness)
\%\textgreater\%\\
filter(Treatment == ``Conv.'') \%\textgreater\%\\
mutate(logRich = log(richness)) \%\textgreater\%\\
head()

\begin{verbatim}

### Grouping and summarising
confusing
``{r} 
microbiome.fungi %>%
  select(SampleID, Crop, Compartment:Fungicide, richness) %>%
  group_by(Treatment, Fungicide) %>%        # define groups
  mutate(logRich = log(richness)) %>%       # uses groups but keeps rows
  summarise(Mean.rich = mean(logRich),      # 1 row per group
            n = n(),
            sd.dev = sd(logRich)) %>%
  mutate(std.err = sd.dev / sqrt(n))

# 2  explore a topic of your choice

Ans: I chose to explore pivot_longer() because reshaping data from wide to long format is often required before plotting and modeling. I wanted to understand how to convert multiple measurement columns into tidy format usable in ggplot and dplyr workflows. 

Example: 


library(tidyr)

long_df <- df %>%
pivot_longer(cols = starts_with("rep"),
names_to = "replicate",
values_to = "value")

It is useful to
•   Easier plotting
•   Easier grouping
•   Tidy data structure


also;
I chose to explore how to plot graphs in R using ggplot2. I picked this topic because plotting is essential for exploring and communicating patterns in data, and I want to become more comfortable creating clear, publication-quality figures. 


library(ggplot2)

ggplot(mtcars, aes(x = wt, y = mpg, color = factor(cyl))) + 
  geom_point(size = 2) +
  geom_smooth(method = "lm", se = TRUE) +
  labs(
    title = "MPG vs Weight by Cylinders",
    x = "Weight (1000 lbs)",
    y = "Miles per gallon",
    color = "Cylinders"
  ) +
  theme_minimal()
\end{verbatim}

I wanted to learn how each piece of the code works: how aes() maps
variables to x, y, and color, how geom\_point() and geom\_smooth() stack
as layers, and how functions like labs() and theme\_minimal() change the
non-data elements of the graph to improve readability.

\section{3 Introduction to renv.}\label{introduction-to-renv.}

• Follow the coding demonstration in class. Turn in notes.

Ans: renv Notes: renv is an R package for project‑specific package
libraries that improves reproducibility and portability. It creates a
local library folder for each project plus a renv.lock file that records
exact package names and versions used.

Purposes: • creates Reproducible R environments • tracks package
versions • allows others to run code without errors

Initialize project:

\begin{Shaded}
\begin{Highlighting}[]
\FunctionTok{install.packages}\NormalTok{(}\StringTok{"renv"}\NormalTok{)}
\end{Highlighting}
\end{Shaded}

\begin{verbatim}
## The following package(s) will be installed:
## - renv [1.2.0]
## These packages will be installed into "~/Reproducibility-GB/BonusQuestion/renv/library/windows/R-4.4/x86_64-w64-mingw32".
## 
## # Installing packages --------------------------------------------------------
## ✔ renv 1.2.0                               [linked from cache]
## Successfully installed 1 package in 36 milliseconds.
\end{verbatim}

\begin{Shaded}
\begin{Highlighting}[]
\NormalTok{renv}\SpecialCharTok{::}\FunctionTok{init}\NormalTok{()}
\FunctionTok{install.packages}\NormalTok{(}\StringTok{"dplyr"}\NormalTok{)}
\end{Highlighting}
\end{Shaded}

\begin{verbatim}
## The following package(s) will be installed:
## - dplyr [1.2.1]
## These packages will be installed into "~/Reproducibility-GB/BonusQuestion/renv/library/windows/R-4.4/x86_64-w64-mingw32".
## 
## # Installing packages --------------------------------------------------------
## ✔ dplyr 1.2.1                              [linked from cache]
## Successfully installed 1 package in 11 milliseconds.
\end{verbatim}

\begin{Shaded}
\begin{Highlighting}[]
\NormalTok{renv}\SpecialCharTok{::}\FunctionTok{snapshot}\NormalTok{()}
\end{Highlighting}
\end{Shaded}

\begin{verbatim}
## - The lockfile is already up to date.
\end{verbatim}

\begin{Shaded}
\begin{Highlighting}[]
\NormalTok{renv}\SpecialCharTok{::}\FunctionTok{restore}\NormalTok{()}
\end{Highlighting}
\end{Shaded}

\begin{verbatim}
## - The library is already synchronized with the lockfile.
\end{verbatim}

Files created:

renv.lock → package versions renv/ folder → project library

It is important because:

ensures reproducibility avoids package conflicts allows project sharing

\end{document}
